%=========== ΛΥΣΗ - 43ο ΘΕΜΑ Δ ===========
\begin{thema}{Δ}
\begin{erwthma}
\item Για τον όγκο και το εμβαδόν της επιφάνειας της σφαίρας ισχύουν
\[
V(t)=\frac{4}{3}\pi R^3(t)
\]
και
\[
E(t)=4\pi R^2(t).
\]
Παραγωγίζοντας τη σχέση του όγκου, παίρνουμε
\[
V'(t)=4\pi R^2(t)R'(t).
\]
Από τη δοσμένη σχέση
\[
V'(t)=-kE(t)
\]
προκύπτει, για κάθε $t\in(0,t_{\mathrm{ολ}})$,
\begin{align*}
4\pi R^2(t)R'(t)
&=-4k\pi R^2(t).
\end{align*}
Επειδή πριν από την ολική τήξη είναι $R(t)>0$, έχουμε
\[
{R'(t)=-k}.
\]
Άρα ο ρυθμός μεταβολής της ακτίνας είναι σταθερός.

Αν
\[
R(0)=R_0,
\]
τότε, αφού $R'(t)=-k$, έχουμε
\[
{R(t)=R_0-kt}.
\]

\item Από τη σχέση
\[
V(2)=\frac{V_0}{8}
\]
έχουμε
\begin{align*}
\frac{4}{3}\pi R^3(2)
&=\frac{1}{8}\cdot\frac{4}{3}\pi R_0^3\\
\Leftrightarrow
R^3(2)&=\frac{R_0^3}{8}\\
\Leftrightarrow
R(2)&=\frac{R_0}{2}.
\end{align*}
Όμως
\[
R(2)=R_0-2k,
\]
οπότε
\begin{align*}
R_0-2k&=\frac{R_0}{2}\\
\Leftrightarrow
R_0&=4k.
\end{align*}

Τη χρονική στιγμή της ολικής τήξης είναι
\[
R(t_{\mathrm{ολ}})=0.
\]
Άρα
\begin{align*}
R_0-kt_{\mathrm{ολ}}&=0\\
\Leftrightarrow
t_{\mathrm{ολ}}&=\frac{R_0}{k}\\
&=4.
\end{align*}
Επομένως
\[
{t_{\mathrm{ολ}}=4\ \text{ώρες}}.
\]

\item Έχουμε
\[
E(t)=4\pi R^2(t).
\]
Παραγωγίζοντας και χρησιμοποιώντας ότι $R'(t)=-k$, παίρνουμε
\begin{align*}
E'(t)
&=8\pi R(t)R'(t)\\
&=-8\pi kR(t).
\end{align*}
Επειδή
\[
R(t)=R_0-kt,
\]
έχουμε
\[
E'(t)=-8\pi k(R_0-kt)
=8\pi k^2t-8\pi kR_0.
\]
Άρα
\[
\left(E'(t)\right)'=8\pi k^2>0.
\]
Επομένως η συνάρτηση $E'$ είναι \textbf{γνησίως αύξουσα} στο διάστημα $(0,t_{\mathrm{ολ}})$.

\item Για $V_0=8\pi$ έχουμε
\[
V(0)=8\pi>1,
\]
ενώ κατά την ολική τήξη
\[
V(t_{\mathrm{ολ}})=0<1.
\]
Η $V$ είναι συνεχής στο $[0,t_{\mathrm{ολ}}]$. Επομένως, από το Θεώρημα Ενδιάμεσων Τιμών, υπάρχει τουλάχιστον μία χρονική στιγμή
\[
t_0\in(0,t_{\mathrm{ολ}})
\]
τέτοια ώστε
\[
V(t_0)=1.
\]

Για κάθε $t\in(0,t_{\mathrm{ολ}})$ είναι
\[
V'(t)=-kE(t)<0,
\]
αφού $k>0$ και $E(t)>0$. Άρα η $V$ είναι γνησίως φθίνουσα στο $[0,t_{\mathrm{ολ}}]$, οπότε η χρονική στιγμή $t_0$ είναι μοναδική.

Συνεπώς υπάρχει ακριβώς μία χρονική στιγμή
\[
{t_0\in(0,t_{\mathrm{ολ}})}
\]
κατά την οποία ο υπολειπόμενος όγκος του πάγου είναι ίσος με $1$.
\end{erwthma}
\end{thema}
%=========== ΤΕΛΟΣ ΛΥΣΗΣ - 43ο ΘΕΜΑ Δ ===========
